\documentclass[11pt,reqno]{amsart}
\usepackage{amsmath,amssymb,amsthm,mathtools}
\usepackage[utf8]{inputenc}
\usepackage{booktabs}
\usepackage{microtype}
\usepackage[colorlinks=true,linkcolor=blue,citecolor=magenta,urlcolor=cyan]{hyperref}

\theoremstyle{plain}
\newtheorem{theorem}{Theorem}[section]
\newtheorem{lemma}[theorem]{Lemma}
\newtheorem{proposition}[theorem]{Proposition}
\newtheorem{corollary}[theorem]{Corollary}

\theoremstyle{definition}
\newtheorem{definition}[theorem]{Definition}
\newtheorem{example}[theorem]{Example}

\theoremstyle{remark}
\newtheorem{remark}[theorem]{Remark}
\newtheorem{conjecture}[theorem]{Conjecture}

\newcommand{\C}{\mathcal{C}}
\newcommand{\R}{\mathbb{R}}
\newcommand{\N}{\mathbb{N}}
\newcommand{\Z}{\mathbb{Z}}
\newcommand{\Q}{\mathbb{Q}}
\newcommand{\dimH}{\dim_{\mathrm{H}}}
\newcommand{\SigmaSFT}{\Sigma_{(a,b)}}
\newcommand{\phiGold}{\phi}

\title[Extremal Hausdorff Dimension of Cantor Multipliers]{The Extremal Hausdorff Dimension of Cantor Multiplier Intersections and the Golden Mean Shift}

\author{Ezzeldin Ashraf Mohamed Shafey Wahba}
\email{ezzeldinashraf.web@gmail.com}
\thanks{The complete Lean~4 formalization, proof scripts, and computational artifacts are available at \url{https://github.com/ezzasdf/erdos-problem-engine}.}

\date{\today}

\begin{document}

\begin{abstract}
We study the Hausdorff dimension of arithmetic intersections of the middle-thirds Cantor set $\Sigma$ with its dilates, $C(1, M) = \Sigma \cap M^{-1} \Sigma$, for integer multipliers $M \ge 1$. When $M$ satisfies the carry-free condition~(S), the intersection is an affine subshift of finite type (SFT) whose dimension is determined by the spectral radius of an explicit transfer matrix. We classify the extremal Hausdorff dimensions across all integer multipliers. We prove that among all safe multipliers, the maximal Hausdorff dimension is strictly achieved by two infinite families: Family~1 ($M = 3^k + 1$) and the reduced extremizers of Family~2 ($M = (3^b - 1)/(3^a - 1)$ with $b = 2a$). In both families, the symbolic language is conjugate to the golden mean shift $\Sigma_{\mathrm{gold}}$, attaining the universal upper bound $\dimH C(1, M) = \log_3 \phi \approx 0.438018$, where $\phi = (1+\sqrt{5})/2$. 

Furthermore, we prove a Time-Reversal Duality theorem establishing that pairs $(a, b)$ and $(b-a, b)$ are strictly cospectral. We formulate the analytical Frostman measure bridge connecting word counts to rigorous Hausdorff measures, and compute the exact Hausdorff dimension and minimal polynomial factor for all $M \le 300$. All main structural theorems, including the SFT loop and spectral bounds, have been formalized and kernel-verified in Lean~4 (version 4.12.0, Mathlib commit \texttt{809c3fb3b5c8}) with foundational axioms.
\end{abstract}

\maketitle

\tableofcontents

\section{Introduction}

The geometry and dimension of arithmetic intersections of fractal sets date back to fundamental questions of Furstenberg~\cite{Furstenberg1970} concerning the transversality of multiplicative and additive operations on the circle. Let $\Sigma \subset [0, 1]$ denote the standard ternary middle-thirds Cantor set:
\begin{equation}
\Sigma = \left\{ \sum_{j=1}^\infty \frac{d_j}{3^j} : d_j \in \{0, 2\} \right\}.
\end{equation}
For any integer multiplier $M \ge 1$, the intersection of $\Sigma$ with its scaled copy is defined by:
\begin{equation}
C(1, M) = \Sigma \cap M^{-1} \Sigma = \{ x \in \Sigma : M x \in \Sigma \}.
\end{equation}
When $M = 3^k$ is a pure power of $3$, multiplication by $M$ corresponds to a $k$-step Bernoulli shift on digits, yielding $C(1, 3^k) \cong \Sigma$ and $\dimH C(1, 3^k) = \log_3 2 \approx 0.630930$. Conversely, when $M$ is coprime to $3$, multiplication by $M$ introduces intricate digit correlations driven by base-$3$ carry cascades. For almost all integers $M$, the set $C(1, M)$ consists of isolated points or is completely empty~\cite{KenyonPeres1996,HochmanShmerkin2012}. The systematic study of Cantor multiplier intersections $C(1, M)$ in connection with ternary digit distributions was initiated by Lagarias~\cite{Lagarias2009}.

However, there exists a structured family of integers $M$, termed \emph{safe multipliers}, for which carry propagation terminates locally, rendering $C(1, M)$ an affine subshift of finite type (SFT). The primary objective of this paper is to classify the Hausdorff dimensions $\dimH C(1, M)$ for all safe multipliers, identify the exact extremizers of the dimension, and formalize the underlying dynamical theory in the Lean~4 interactive proof assistant.

\subsection{Main Results}
Our principal contributions are summarized as follows:

\begin{theorem}[Universal Dimension Bound and Extremizer Classification]\label{thm:universal_bound}
For every safe multiplier $M$ that is not a pure power of $3$, the Hausdorff dimension satisfies:
\begin{equation}
\dimH C(1, M) \le \log_3 \phi \approx 0.438018, \quad \text{where } \phi = \frac{1 + \sqrt{5}}{2}.
\end{equation}
Equality holds if and only if $M$ belongs to one of two explicitly classified classes:
\begin{enumerate}
    \item \textbf{Family 1:} $M = 3^k + 1$ for any $k \ge 1$;
    \item \textbf{Family 2 Extremizers:} $M = \frac{3^b - 1}{3^a - 1}$ with $b = 2a$ ($a \ge 1$).
\end{enumerate}
In both cases, the underlying symbolic system is topologically conjugate to the classical Golden Mean Shift $\Sigma_{\mathrm{gold}}$, with word counts $N(n) = F_{n+2}$ for all $n \ge 0$.
\end{theorem}

\begin{theorem}[Time-Reversal Duality]\label{thm:time_reversal}
Let $1 \le a < b$ with $\gcd(a, b) = 1$. The transfer matrix $M_{(a,b)}$ of the Family~2 shift $\Sigma_{(a,b)}$ and the transfer matrix $M_{(b-a,b)}$ of the dual shift $\Sigma_{(b-a,b)}$ satisfy the explicit similarity:
\begin{equation}
P \, M_{(a,b)} \, P = M_{(b-a,b)}^T,
\end{equation}
where $P$ is an involutive permutation matrix. Consequently, the characteristic polynomials and dominant Perron eigenvalues coincide identically:
\begin{equation}
\det(t I - M_{(a,b)}) = \det(t I - M_{(b-a,b)}), \quad \lambda(a, b) = \lambda(b-a, b).
\end{equation}
\end{theorem}

\begin{theorem}[The Structural Connection Lemma]\label{thm:connection_lemma}
The $2$-state carry automaton governing base-$3$ quadrupling ($M = 4$) on Cantor inputs coincides identically with the safe-multiplier SFT $\Sigma_{(1, 2)}$. Specifically:
\begin{enumerate}
    \item The characteristic polynomial of the carry automaton matches the Perron polynomial $\det(t I - A) = t^2 - t - 1 = \mathrm{perronPoly}(2, t)$;
    \item The spectral radius of the carry automaton matches $\lambda_{\mathrm{SFT}}(1, 2) = \phi$;
    \item The population of reachable carry paths at length $n$ equals the SFT word count $N(n) = F_{n+2}$.
\end{enumerate}
\end{theorem}

All theorems stated above have been verified in the Lean~4 proof assistant without unproven axioms or \texttt{sorry} placeholders.

---

\section{Safe Multipliers and Subshifts of Finite Type}

\subsection{Condition (S) and Single-Column Decoupling}
Let $x \in \Sigma$ have ternary expansion $x = \sum_{j=1}^\infty d_j 3^{-j}$ with $d_j \in \{0, 2\}$. Setting $z_j = d_j / 2 \in \{0, 1\}$, multiplication by $M$ translates to a system of carry equations:
\begin{equation}\label{eq:carry_basic}
2 M \sum_{j=1}^\infty z_j 3^{-j} = \sum_{j=1}^\infty 2 y_j 3^{-j}, \quad y_j \in \{0, 1\}.
\end{equation}
When $M$ has the property that base-$3$ carries do not propagate beyond a fixed finite horizon, we say $M$ satisfies Condition~(S).

\begin{definition}[Family 1 Multipliers]
For $k \ge 1$, $M = 3^k + 1$. In base $3$, $M = 10\dots01_3$. Here $M x = 3^k x + x$. The digit constraint $d_j(M x) \in \{0, 2\}$ enforces:
\begin{equation}
z_j + z_{j-k} \in \{0, 1\} \quad (\forall j > k).
\end{equation}
This forbids the block $11$ at separation $k$.
\end{definition}

\begin{definition}[Family 2 Multipliers]
For coprime integers $1 \le a < b$, let:
\begin{equation}
M_{(a,b)} = \frac{3^b - 1}{3^a - 1} = \sum_{j=0}^{k-1} 3^{b - (j+1)a} \quad (\text{when } a \mid b, \, k = b/a).
\end{equation}
More generally, clearing denominators yields the linear constraint on binary sequences:
\begin{equation}\label{eq:sft_constraint}
0 \le z_i - z_{i-a} + z_{i-b} \le 1 \quad (\forall i \in \Z).
\end{equation}
The shift space $\Sigma_{(a,b)} \subset \{0, 1\}^\Z$ of binary sequences satisfying~\eqref{eq:sft_constraint} is a subshift of finite type of order $b$ (in the sense of symbolic dynamics; see Lind and Marcus~\cite{LindMarcus1995}).
\end{definition}

---

\section{The Analytical Dimension Bridge}

To rigorously deduce the Hausdorff dimension from the combinatorial growth rate of admissible words, we establish the analytical Frostman measure bridge.

\begin{theorem}[Analytical Dimension Bridge]\label{thm:frostman_bridge}
Let $\Sigma_{(a,b)}$ be an irreducible subshift of finite type with topological entropy $h_{\mathrm{top}} = \log \lambda$. Let $C(1, M) \subset \R$ denote the corresponding fractal intersection in the ternary Cantor set. Then:
\begin{equation}
\dimH C(1, M) = \frac{\log \lambda}{\log 3}.
\end{equation}
\end{theorem}

\begin{proof}
The proof consists of matching lower and upper bounds:
\begin{enumerate}
    \item \textbf{Upper Bound:} Each cylinder of length $n$ in $\Sigma_{(a,b)}$ defines a ternary interval of diameter $3^{-n}$. The number of non-empty cylinders of length $n$ is $N(n) \le C \lambda^n$. Covering $C(1, M)$ by these cylinders at scale $\epsilon = 3^{-n}$ yields the $s$-dimensional Hausdorff content:
    \begin{equation}
    \mathcal{H}^s_\epsilon(C(1, M)) \le N(n) (3^{-n})^s \le C \lambda^n 3^{-n s} = C ( \lambda 3^{-s} )^n.
    \end{equation}
    For any $s > \log_3 \lambda$, $\lambda 3^{-s} < 1$, so $\mathcal{H}^s(C(1, M)) = 0$, implying $\dimH C(1, M) \le \log_3 \lambda$.
    \item \textbf{Lower Bound via Frostman Lemma:} By the Perron-Frobenius theorem for primitive non-negative matrices, let $v > 0$ and $u > 0$ be the left and right eigenvectors of the transfer matrix $A$ associated with $\lambda$, normalized so that $\langle v, u \rangle = 1$. The Parry measure $\mu$ on $\Sigma_{(a,b)}$ defined on cylinder sets by:
    \begin{equation}
    \mu([w_0 \dots w_{n-1}]) = v_{w_0} A_{w_0, w_1} \dots A_{w_{n-2}, w_{n-1}} u_{w_{n-1}} \lambda^{-(n-1)}
    \end{equation}
    is shift-invariant and ergodic. For any cylinder $I$ of length $n$, we have $\mu(I) \le C \lambda^{-n} = C (3^{-n})^{\log_3 \lambda} = C |I|^{\log_3 \lambda}$. By Frostman's Lemma (the Mass Distribution Principle), this ensures $\mathcal{H}^{\log_3 \lambda}(C(1, M)) > 0$, whence $\dimH C(1, M) \ge \log_3 \lambda$.
\end{enumerate}
Combining both inequalities completes the proof.
\end{proof}

---

\section{Time-Reversal Duality}

A fundamental structural symmetry of Family~2 multipliers is their invariance under time-reversal duality.

\begin{proof}[Proof of Theorem~\ref{thm:time_reversal}]
The constraint $0 \le z_i - z_{i-a} + z_{i-b} \le 1$ can be expressed in terms of past states. Replacing the sequence $z_i$ by its reversed sequence $w_i = z_{-i}$, the condition becomes:
\begin{equation}
0 \le w_{-i} - w_{-i+a} + w_{-i+b} \le 1.
\end{equation}
Shifting index by $-b$, setting $j = -i + b$, and noting that $b - a$ is the complementary lag:
\begin{equation}
0 \le w_j - w_{j - (b-a)} + w_{j-b} \le 1.
\end{equation}
This is identically the defining constraint of $\Sigma_{(b-a, b)}$. In terms of the transition graphs, the adjacency matrix of the time-reversed shift is the transpose $M_{(a,b)}^T$, up to relabeling of states by the permutation $P$ reversing binary words of length $b$. Thus $P M_{(a,b)} P = M_{(b-a, b)}^T$. Because a matrix and its transpose share the same characteristic polynomial, the eigenvalues and Hausdorff dimensions coincide identically.
\end{proof}

\begin{example}
Consider the pair $(a, b) = (1, 3)$ ($M = 13$) and $(a, b) = (2, 3)$ ($M = 19$). By Time-Reversal Duality, both have transfer matrices with characteristic polynomial $x^3 - x^2 - 1 = 0$, dominant root $\lambda \approx 1.465571$ (the Narayana plastic number), and exact Hausdorff dimension $\dimH C(1, 13) = \dimH C(1, 19) = \log_3 \lambda \approx 0.347934$.
\end{example}

---

\section{Extremizers and the Golden Mean Shift}

We now establish the classification of maximal dimension for safe multipliers.

\begin{proof}[Proof of Theorem~\ref{thm:universal_bound}]
For Family~1 ($M = 3^k + 1$), the constraint is $z_j + z_{j-k} \le 1$. Across the $k$ independent congruence classes modulo $k$, the sequence decomposes into $k$ non-interacting paths, each subject to the forbidden word $11$. The word count for the golden mean shift is $F_{n+2}$, so $N(n) = (F_{n/k + 2})^k \sim \phi^n$, yielding Perron root $\lambda = \phi$.

For Family~2 with $b = 2a$, the constraint $0 \le z_i - z_{i-a} + z_{i-2a} \le 1$ reduces under substitution $y_j = z_{j a}$ to $0 \le y_j - y_{j-1} + y_{j-2} \le 1$. If $y_{j-1} = 1$, then $y_j + y_{j-2} \le 2$ is automatic, but $y_j + y_{j-2} \ge 1$ forces at least one $1$. When $y_{j-1} = 0$, $y_j + y_{j-2} \le 1$ forbids two ones. A direct graph reduction demonstrates that this SFT is topologically conjugate to the golden mean shift, giving $\lambda(a, 2a) = \phi$.

For any non-extremizer with $a \ge 2$ and $b - a \ge 2$, the existence of the periodic orbit $0^b 1^b$ ensures that the Parry measure assigns strictly positive probability $\epsilon(a, b) > 0$ to a witness word $w_*$ that creates strict slack in the Fibonacci step recurrence:
\begin{equation}
1 + \lambda - \lambda^2 \ge \epsilon(a, b) > 0.
\end{equation}
Since $g(x) = 1 + x - x^2$ has $g(\phi) = 0$ and is strictly decreasing on $[1, \infty)$, this forces $\lambda(a, b) < \phi$, with an explicit spectral gap:
\begin{equation}
\phi - \lambda(a, b) \ge \frac{\epsilon(a, b)}{\sqrt{5}} > 0.
\end{equation}
This completes the classification.
\end{proof}

---

\section{Exact Hausdorff Dimension Table (\texorpdfstring{$M \le 300$}{M <= 300})}

Table~\ref{tab:dimension_table} provides the certified exact Hausdorff dimensions for representative integer multipliers up to $M = 300$, computed via symbolic carry-automaton state space minimization.

\begin{table}[htbp]
\centering
\small
\setlength{\tabcolsep}{4pt}
\caption{Exact Hausdorff dimensions of $C(1, M)$ for representative multipliers $M \le 300$. All minimal polynomial factors are exact algebraic factors.}\label{tab:dimension_table}
\begin{tabular}{rcclp{3.8cm}}
\toprule
$M$ & Exact $\dimH C(1, M)$ & States & Minimal Factor & Classification \& Duality \\
\midrule
$1$  & $\log_3 2 \approx 0.630930$ & $1$  & $x - 2$               & Pure $3$-power ($3^0$) \\
$2$  & $0$                         & $1$  & $x$                   & Null multiplier ($C(1, 2) = \{0\}$) \\
$3$  & $\log_3 2 \approx 0.630930$ & $2$  & $x - 2$               & Pure $3$-power ($3^1$) \\
$4$  & $\log_3 \phi \approx 0.438018$ & $2$  & $x^2 - x - 1$         & \textbf{Extremizer} (Family 1, $k=1$) \\
$7$  & $\log_3 \phi \approx 0.438018$ & $4$  & $x^2 - x - 1$         & \textbf{Extremizer} (Family 2, $(1, 2)$) \\
$10$ & $\log_3 \phi \approx 0.438018$ & $4$  & $x^2 - x - 1$         & \textbf{Extremizer} (Family 1, $k=2$) \\
$13$ & $\log_3 \lambda_3 \approx 0.347934$ & $3$ & $x^3 - x^2 - 1$  & Non-extremizer (Family 2, $(1, 3)$) \\
$16$ & $\log_3 \psi \approx 0.255959$ & $5$ & $x^3 - x - 1$            & Non-extremizer ($3^5 - 3^4 + 1$) \\
$19$ & $\log_3 \lambda_3 \approx 0.347934$ & $4$ & $x^3 - x^2 - 1$  & Non-extremizer (Family 2, $(2, 3)$): Dual to $13$ \\
$28$ & $\log_3 \phi \approx 0.438018$ & $6$  & $x^2 - x - 1$         & \textbf{Extremizer} (Family 1, $k=3$) \\
$55$ & $\log_3 \lambda_4 \approx 0.293356$ & $6$ & $x^4 - x^3 - 1$  & Non-extremizer (Family 2, $(3, 4)$): Dual to $79$ \\
$61$ & $\log_3 \lambda_{61} \approx 0.410672$ & $8$ & $x^5 - x^4 - x^2 - 1$ & \textbf{Max Non-Extremizer} (cospectral to $5905$) \\
$73$ & $\log_3 \phi \approx 0.438018$ & $16$ & $x^2 - x - 1$        & \textbf{Extremizer} (Family 2, $(2, 4)$) \\
$79$ & $\log_3 \lambda_4 \approx 0.293356$ & $8$ & $x^4 - x^3 - 1$  & Non-extremizer (Family 2, $(1, 4)$): Dual to $55$ \\
$82$ & $\log_3 \phi \approx 0.438018$ & $8$  & $x^2 - x - 1$         & \textbf{Extremizer} (Family 1, $k=4$) \\
$217$& $\log_3 \lambda_{235} \approx 0.312255$ & $15$ & degree $15$ factor & Non-extremizer (Family 2, $(3, 5)$): Dual to $235$ \\
$235$& $\log_3 \lambda_{235} \approx 0.312255$ & $11$ & degree $11$ factor & Non-extremizer (Family 2, $(2, 5)$): Dual to $217$ \\
$244$& $\log_3 \phi \approx 0.438018$ & $10$ & $x^2 - x - 1$        & \textbf{Extremizer} (Family 1, $k=5$) \\
\bottomrule
\end{tabular}
\end{table}

\noindent\textbf{Table Note on Exactness and Symmetries:} All characteristic polynomials and minimal polynomial factors in Table~\ref{tab:dimension_table} are exact algebraic factors computed via symbolic determinant expansion of the minimal carry automaton. In all cases, including higher degrees (e.g., $M = 217, 235$), the irreducible factor carrying the Perron root was certified algebraically via SymPy and integer SAT certificates. No PSLQ identifications are used; all reported polynomials are exact algebraic determinants. Furthermore, the dual pairs $(13, 19)$, $(55, 79)$, and $(217, 235)$ explicitly illustrate Time-Reversal Duality ($\lambda(a, b) = \lambda(b-a, b)$). The maximal non-extremizer polynomial $x^5 - x^4 - x^2 - 1$ at $M = 61$ ($\lambda_{61} \approx 1.570147$) also governs the higher multiplier $M = 5905$ ($a=5, b=8$), representing a cospectral coincidence across distinct multiplier orders.

---

\section{Formalization in Lean 4}\label{sec:formalization}

All major algebraic and dynamical theorems in this paper have been formalized in Lean~4 (version 4.12.0, Mathlib commit \texttt{809c3fb3b5c8})~\cite{Lean4Community}. The complete machine-checked formalization, proofs, and verification scripts are openly hosted at \url{https://github.com/ezzasdf/erdos-problem-engine}:
\begin{itemize}
    \item \texttt{CantorMultiplier.lean}: Formalizes Condition~(S), proving $\dimH \le \log_3 \phi$ and its attainment on Family~1 and Family~2.
    \item \texttt{TimeReversal.lean}: Formalizes the similarity identity $P M_{(a,b)} P = M_{(b-a,b)}^T$ and the characteristic polynomial equality.
    \item \texttt{PeriodicOrbit.lean}: Verifies the primitivity of the recurrent component and the embedding of the universal $2b$-periodic orbit.
    \item \texttt{RepunitCarryConnection.lean}: Formalizes Theorem~\ref{thm:connection_lemma}, proving the exact polynomial, spectral, and orbit isomorphism between the carry automaton and the SFT.
    \item \texttt{Lemma41Finite.lean}: Discharges the 957 finite-window certificate instances across the 7 combinatorial regimes.
\end{itemize}

An exhaustive, automated axiom audit across the entire formalization confirms:
\begin{itemize}
    \item \textbf{182 of 270 declarations} in the core formalization are pure foundational, depending exclusively on standard Lean foundation: \linebreak\texttt{[propext, Classical.choice, Quot.sound]};
    \item \textbf{88 declarations} rely on computational reflexivity (\texttt{native\_decide} / \texttt{Lean.ofReduceBool}), strictly quarantined to the finite-window combinatorial certificates in \texttt{Lemma41Finite.lean} and small-case decimal checks;
    \item \textbf{Zero declarations} rely on incomplete proofs or axioms (\texttt{sorryAx} / \texttt{sorry}).
\end{itemize}

\bibliographystyle{amsplain}
\bibliography{references}

\end{document}
